\chapter{周秦卷索引}

本索引只收录本卷已经写入正文、且可由链接抵达的事件，以及已经置入正文的图件；它不是事件账本的替代品。年代有争议的节点仍依正文的说明处理，未写成正文的账本节点不在此增列。

\section{中国事件索引}

\subsection*{西周}
\begin{itemize}
  \item \eventref{WZ-1046-MUYE}{西周·牧野}；\eventref{WZ-1040-DONGZHENG}{西周·三监东征}；\eventref{WZ-EAST-CHENGZHOU}{西周·营建成周}。
  \item \eventref{WZ-CHENGKANG-CONSOLIDATION}{西周·成康巩固}；\eventref{WZ-MU-TRAVELS}{西周·穆王远行}；\eventref{WZ-0977-ZHAONAN}{西周·昭王南征}。
  \item \eventref{WZ-0841-GONGHE}{西周·共和}；\eventref{WZ-0828-XUAN-ACCESSION}{西周·宣王即位}；\eventref{WZ-XUAN-REVIVAL}{西周·宣王中兴}；\eventref{WZ-0771-HAOJING}{西周·镐京陷落}。
\end{itemize}

\subsection*{春秋}
\begin{itemize}
  \item \eventref{SA-0770-EAST}{春秋·东迁}；\eventref{SA-0770-QIN-ENFEOFF}{春秋·秦获封关中}；\eventref{SA-0722-CHUNQIU}{春秋·鲁隐公元年}；\eventref{SA-0707-XUGE}{春秋·繻葛之战}；\eventref{SA-0704-CHU-KING}{春秋·楚君称王}。
  \item \eventref{SA-0685-QIHUAN}{春秋·齐桓公即位}；\eventref{SA-0656-SHAOLING}{春秋·召陵之盟}；\eventref{SA-0651-KUIQIU}{春秋·葵丘之会}；\eventref{SA-0636-JIN-WEN}{春秋·晋文公即位}；\eventref{SA-0632-CHENGPU}{春秋·城濮之战}；\eventref{SA-0627-XIAO}{春秋·殽之战}。
  \item \eventref{SA-0597-BI}{春秋·邲之战}；\eventref{SA-0575-YANLING}{春秋·鄢陵之战}；\eventref{SA-0546-MIBING}{春秋·弭兵之会}；\eventref{SA-0506-BOJU}{春秋·柏举之战}；\eventref{SA-0494-FUJIAO}{春秋·夫椒之战}。
  \item \eventref{SA-0484-WUZIXU-DEATH}{春秋·伍子胥之死}；\eventref{SA-0482-HUANGCHI}{春秋·黄池之会}；\eventref{SA-0481-LIN}{春秋·获麟与陈氏政变}；\eventref{SA-0481-END}{春秋·记事终点}；\eventref{SA-0473-WU-FALL}{春秋·越灭吴}。
\end{itemize}

\subsection*{战国与秦}
\begin{itemize}
  \item \eventref{WQ-0453-ZHI}{战国·三家灭智}；\eventref{WQ-0408-WEI}{战国·魏国早期改革}；\eventref{WQ-0403-THREEJIN}{战国·三家分晋}；\eventref{WQ-0386-TIANQI}{战国·田氏代齐}；\eventref{WQ-0381-WUQI-CHU}{战国·吴起变法}；\eventref{WQ-0375-HAN}{战国·韩灭郑}。
  \item \eventref{WQ-0362-WEIQIN}{战国·少梁之战}；\eventref{WQ-0356-SHANGYANG}{战国·商鞅变法}；\eventref{WQ-0354-SHENBUHAI}{战国·申不害相韩}；\eventref{WQ-0354-GUILIN}{战国·桂陵之战}；\eventref{WQ-0342-MALING}{战国·马陵之战}；\eventref{WQ-0334-KINGS}{战国·齐魏称王}。
  \item \eventref{WQ-0317-YAN}{战国·子之之乱}；\eventref{WQ-0314-QI-YAN}{战国·齐伐燕}；\eventref{WQ-0307-HUFU}{战国·胡服骑射}；\eventref{WQ-0307-QIN-SUCCESSION}{秦·昭襄王即位}；\eventref{WQ-0295-SHAQIU}{战国·沙丘之乱}；\eventref{WQ-0293-YIQUE}{战国·伊阙之战}。
  \item \eventref{WQ-0284-QI}{战国·燕破齐}；\eventref{WQ-0279-TIANDAN}{战国·田单复齐}；\eventref{WQ-0266-FANJU}{秦·范雎入秦与远交近攻}；\eventref{WQ-0262-SHANGDANG}{战国·上党归赵}；\eventref{WQ-0260-CHANGPING}{战国·长平之战}。
  \item \eventref{WQ-0256-ZHOUDIE}{战国·东周亡}；\eventref{WQ-0249-EASTZHOU}{战国·东周君国亡}；\eventref{QIN-0230-UNIFICATION}{秦·统一六国}；\eventref{QIN-0226-YAN}{秦·破蓟}；\eventref{QIN-0222-YAN}{秦·灭燕}。
  \item \eventref{QIN-0213-EMPIRE}{秦·高强度治理}；\eventref{QIN-0210-SUCCESSION}{秦·沙丘继承}；\eventref{QIN-0209-FALL}{秦·起事与覆亡}、\eventref{QIN-0209-CHENSHENG}{秦·大泽起事}；\eventref{QIN-0208-ZHANGHAN}{秦·章邯出关}；\eventref{QIN-0207-LISI}{秦·李斯被杀}；\eventref{QIN-0207-ZIYING}{秦·子婴即位}；\eventref{QIN-0207-JULU}{秦·巨鹿之战}；\eventref{QIN-0206-SURRENDER}{秦·子婴降汉}。
\end{itemize}

\section{政体与国家索引}

\begin{description}
  \item[周王室与西周] 牧野后的王室建构、东方经营与晚期危机：\eventref{WZ-1046-MUYE}{西周·牧野}、\eventref{WZ-EAST-CHENGZHOU}{西周·营建成周}、\eventref{WZ-0841-GONGHE}{西周·共和}、\eventref{WZ-0771-HAOJING}{西周·镐京陷落}。
  \item[东周王室] 东迁以后保有册命与名分，至战国退出政治舞台：\eventref{SA-0770-EAST}{春秋·东迁}、\eventref{SA-0707-XUGE}{春秋·繻葛之战}、\eventref{WQ-0403-THREEJIN}{战国·三家分晋}、\eventref{WQ-0256-ZHOUDIE}{战国·东周亡}、\eventref{WQ-0249-EASTZHOU}{战国·东周君国亡}。
  \item[齐] 齐桓与会盟、田氏易姓、伐燕及复齐：\eventref{SA-0685-QIHUAN}{春秋·齐桓公即位}、\eventref{SA-0656-SHAOLING}{春秋·召陵之盟}、\eventref{WQ-0386-TIANQI}{战国·田氏代齐}、\eventref{WQ-0314-QI-YAN}{战国·齐伐燕}、\eventref{WQ-0284-QI}{战国·燕破齐}、\eventref{WQ-0279-TIANDAN}{战国·田单复齐}。
  \item[晋、赵、魏、韩] 晋的霸权与卿族裂解，及其后的三国竞争：\eventref{SA-0636-JIN-WEN}{春秋·晋文公即位}、\eventref{SA-0632-CHENGPU}{春秋·城濮之战}、\eventref{SA-0575-YANLING}{春秋·鄢陵之战}、\eventref{WQ-0453-ZHI}{战国·三家灭智}、\eventref{WQ-0403-THREEJIN}{战国·三家分晋}；另见\eventref{WQ-0408-WEI}{战国·魏国早期改革}、\eventref{WQ-0375-HAN}{战国·韩灭郑}、\eventref{WQ-0307-HUFU}{战国·胡服骑射}与\eventref{WQ-0260-CHANGPING}{战国·长平之战}。
  \item[楚] 从自称王到江汉大国的北上与制度压力：\eventref{SA-0704-CHU-KING}{春秋·楚君称王}、\eventref{SA-0597-BI}{春秋·邲之战}、\eventref{SA-0506-BOJU}{春秋·柏举之战}、\eventref{WQ-0381-WUQI-CHU}{战国·吴起变法}。
  \item[秦] 关中边缘、变法与统一帝国：\eventref{SA-0770-QIN-ENFEOFF}{春秋·秦获封关中}、\eventref{SA-0627-XIAO}{春秋·殽之战}、\eventref{WQ-0356-SHANGYANG}{战国·商鞅变法}、\eventref{WQ-0350-QINREFORM}{战国·秦迁咸阳与第二次变法}、\eventref{QIN-0230-UNIFICATION}{秦·统一六国}、\eventref{QIN-0206-SURRENDER}{秦·子婴降汉}。
  \item[鲁、郑、宋] 编年书写、中原要冲与弭兵调停：\eventref{SA-0722-CHUNQIU}{春秋·鲁隐公元年}、\eventref{SA-0707-XUGE}{春秋·繻葛之战}、\eventref{SA-0546-MIBING}{春秋·弭兵之会}、\eventref{SA-0481-LIN}{春秋·获麟与陈氏政变}。
  \item[吴、越] 下游水网中的北上争盟与相互吞并：\eventref{SA-0506-BOJU}{春秋·柏举之战}、\eventref{SA-0494-FUJIAO}{春秋·夫椒之战}、\eventref{SA-0482-HUANGCHI}{春秋·黄池之会}、\eventref{SA-0473-WU-FALL}{春秋·越灭吴}。
  \item[燕] 内乱、复仇式动员与秦军东进：\eventref{WQ-0317-YAN}{战国·子之之乱}、\eventref{WQ-YAN-ZHAO}{战国·燕昭王蓄力}、\eventref{WQ-0284-QI}{战国·燕破齐}、\eventref{QIN-0226-YAN}{秦·破蓟}、\eventref{QIN-0222-YAN}{秦·灭燕}。
\end{description}

\section{主题索引}

\begin{description}
  \item[战争] 从王朝奠基的\eventref{WZ-1046-MUYE}{西周·牧野}，到春秋的\eventref{SA-0632-CHENGPU}{春秋·城濮之战}、\eventref{SA-0597-BI}{春秋·邲之战}、\eventref{SA-0506-BOJU}{春秋·柏举之战}，再到战国的\eventref{WQ-0354-GUILIN}{战国·桂陵之战}、\eventref{WQ-0342-MALING}{战国·马陵之战}、\eventref{WQ-0293-YIQUE}{战国·伊阙之战}与\eventref{WQ-0260-CHANGPING}{战国·长平之战}。参见秦末的\eventref{QIN-0207-JULU}{秦·巨鹿之战}。
  \item[改革与国家能力] \eventref{WQ-0408-WEI}{战国·魏国早期改革}、\eventref{WQ-WEI-WUQI}{战国·魏武卒}、\eventref{WQ-0381-WUQI-CHU}{战国·吴起变法}、\eventref{WQ-0354-SHENBUHAI}{战国·申不害相韩}、\eventref{WQ-0356-SHANGYANG}{战国·商鞅变法}与\eventref{WQ-0350-QINREFORM}{战国·秦迁咸阳与第二次变法}。另见图“战国变法关系示意”。
  \item[经济、动员与交通] 西周的成周和东方网络以\eventref{WZ-EAST-CHENGZHOU}{西周·营建成周}为入口；战国则可从\eventref{WQ-0350-QINREFORM}{战国·秦迁咸阳与第二次变法}、\eventref{WQ-0262-SHANGDANG}{战国·上党归赵}与\eventref{WQ-0260-CHANGPING}{战国·长平之战}观察田界、县邑、道路与补给怎样进入竞争。秦帝国的日常征发见\eventref{QIN-0213-EMPIRE}{秦·高强度治理}。
  \item[文化、礼制与书写] 西周王权的礼制语言见\eventref{WZ-CHENGKANG-CONSOLIDATION}{西周·成康巩固}；春秋编年视角的起点和终点见\eventref{SA-0722-CHUNQIU}{春秋·鲁隐公元年}、\eventref{SA-0481-END}{春秋·记事终点}。齐的稷下、游说与法家讨论置于战国“文化与思想”一节，并与\eventref{WQ-0386-TIANQI}{战国·田氏代齐}所后的齐国路径并读。
  \item[边疆、地理与迁徙] 南向通道的限度见\eventref{WZ-0977-ZHAONAN}{西周·昭王南征}，西北压力见\eventref{WZ-XUAN-REVIVAL}{西周·宣王中兴}；秦在关中的位置见\eventref{SA-0770-QIN-ENFEOFF}{春秋·秦获封关中}；赵的北方应变见\eventref{WQ-0307-HUFU}{战国·胡服骑射}；燕的辽东退守与覆亡见\eventref{QIN-0226-YAN}{秦·破蓟}、\eventref{QIN-0222-YAN}{秦·灭燕}。
  \item[史料与争议] 年代与材料层次的提示集中见\eventref{WZ-1046-MUYE}{西周·牧野}、\eventref{WZ-MU-TRAVELS}{西周·穆王远行}、\eventref{WZ-0841-GONGHE}{西周·共和}、\eventref{SA-0481-END}{春秋·记事终点}、\eventref{WQ-0279-TIANDAN}{战国·田单复齐}与\eventref{QIN-0209-CHENSHENG}{秦·大泽起事}。相关正文区分金文、传世编年与后出叙事的证据范围，不把后世故事当作现场记录。
\end{description}

\section{图表与地图索引}

\begin{description}
  \item[西周] “时间线：从牧野到东迁”；“空间总览：王畿、东方与边疆”；“牧野战事空间示意”；“宗周、成周与东方网络示意”。
  \item[春秋] “空间总览：诸侯怎样围绕中原展开”；“繻葛之战空间示意”；“城濮战役空间示意”；“柏举战役空间示意”。
  \item[战国与秦] “空间总览：七国竞争与秦的东进通道”；“战国变法关系示意”；“长平战场空间示意”；“燕伐齐空间示意”；“秦统一六国空间示意”。
\end{description}
